\documentclass{article}
\usepackage[margin=1in]{geometry}

\setcounter{secnumdepth}{3}
% Useful packages
\usepackage{graphicx}
\usepackage{amsmath}
\usepackage{amssymb}
\usepackage{hyperref}
\usepackage{braket}
\usepackage{dsfont}
%\usepackage{showlabels}
\usepackage{amsthm}
\newtheorem{theorem}{Theorem}
\newtheorem{lemma}[theorem]{Lemma}
\newtheorem{corollary}[theorem]{Corollary}
\newtheorem{proposition}[theorem]{Proposition}
\theoremstyle{definition}
\newtheorem{definition}{Definition}[section]
\newtheorem{prop}{Conjecture}
\newtheorem{cor}{Corollary}
%-------------------packages-------------------------
\usepackage{enumerate}
\usepackage[utf8]{inputenc}
\usepackage{amsfonts}
\usepackage{amsbsy, mathtools}
% \usepackage{braket} % already loaded above
% \usepackage{graphicx} % already loaded above
\usepackage{xcolor}
%\usepackage{multirow}
% \usepackage{hyperref} % already loaded above
%\usepackage{caption}
%\usepackage{subcaption}
%\captionsetup{compatibility=false}
\usepackage{natbib}
%\usepackage{dcolumn}% Align table columns on decimal point
\usepackage{bm}% bold math
\usepackage{soul}
\usepackage{showlabels}

%Includes "References" in the table of contents
\usepackage[nottoc]{tocbibind}

\usepackage{bbm}
%\newtheorem{theorem}{Theorem}
%\newtheorem{lemma}[theorem]{Lemma}
%\newtheorem{corollary}[theorem]{Corollary}
%\theoremstyle{definition}
%\newtheorem{definition}{Definition}[section]
\usepackage{tikz}
\usepackage{quantikz}
\usepackage{textcomp}
\usepackage{caption}
\usepackage{microtype}
\newcommand{\Tr}{\text{tr}}
\newcommand{\mA}{\mathcal{A}}
\newcommand{\mB}{\mathcal{B}}
\newcommand{\nn}{\nonumber}
\newcommand{\mL}{\mathcal{L}}
\newcommand{\lb}{\left(}
\newcommand{\rb}{\right)}
\newcommand{\mU}{\mathcal{U}}
\newcommand{\mH}{\mathcal{H}}
\usepackage[T1]{fontenc}
\usepackage{lmodern}
\usepackage{subcaption}
\newcommand{\tOO}{\tilde{\mathcal{O}}}
\newcommand{\mO}{\mathcal{O}}
\newcommand{\mK}{\mathcal{K} }
\newcommand{\ep}{\epsilon}
\newcommand{\mbR}{\mathbb{R}}
\newcommand{\mM}{\mathcal{M}}
\newcommand{\mT}{\mathcal{T}}
\newcommand{\NL}[1]{\textcolor{blue}{#1}}
\newcommand{\p}{\partial}
\newcommand{\tr}{\text{tr}}
\newcommand{\tT}{\tilde{T}}

\usepackage{graphicx} % Required for inserting images

\title{Decay of correlations in ensemble-averaging theories}
\author{Nima Lashkari}
\date{January 2026}

\begin{document}
    \maketitle
    \tableofcontents

    \section*{Results}

\begin{enumerate}
\item \textbf{Time-independent infinite-temperature ensemble.}
For time-independent ensemble-averaged quantum theories with a fixed-$q$
large-$N$ limit, we find that the $k$-twirl map has the leading form
\begin{equation}
    \mathbb E_g\!\left[U^{\otimes k}\otimes(U^*)^{\otimes k}\right]
    \simeq e^{-t^2\widehat G_{\rm eff}^{(k)}}.
\end{equation}
We call this a Gaussian decay.  The symbol $\simeq$ denotes equality to
leading order in the fixed-$q$, fixed-$k$ large-$N$ expansion in the
diffusive regime described below.

\item \textbf{Brownian ensemble.}
For Brownian random couplings, the disorder average can be performed exactly
at finite $N$, and
\begin{equation}
    \mathbb E_g\!\left[U^{\otimes k}\otimes(U^*)^{\otimes k}\right]
    =e^{-T\widehat G_{\rm eff}^{(k)}}.
\end{equation}
The algebraic form of the generator is the same as in the time-independent
problem, although the static variance and the Brownian noise strength have
different dimensions.

\item We study these generators in explicit fermionic SYK and bosonic Spin
SYK systems.
\end{enumerate}


\section*{To do list:}

\begin{enumerate}
\item Generalize the derivation of the $k$-twirl map in the time-independent
case at large $N$ to finite temperature.

\item Understand the physical implications of the $t^2$ in the exponent of
the $k$-twirl map in the large-$N$ limit of regular SYK.  At late time the
decay seems too fast.

\item Analyze the spectrum of $G_{\rm eff}^{(k)}$, the invariant subalgebra,
and establish the relevant spectral gap.

\item Think about the gluing assumption 1.4 in Nick's notes.  Does it hold in
SYK?  Relatedly, if we couple two Brownian SYK models as in Swingle's 2021
paper, does the resulting protocol resemble the Alice--Bob construction in
which one applies local $k$-twirls on $A$ and $B$, swaps $r$ qubits, and
applies the local $k$-twirls again?

\item The Lindblad generator is a sum over a hypergraph with hyperedges
indexed by distinct sets $\{i_1,\ldots,i_q\}$.  Generalize the corresponding
local-qubit constructions to hypergraphs with a fermion on each vertex.
Does the gluing lemma continue to hold for fermionic degrees of freedom?

\item Understand the double-scaled limit $N,q\to\infty$ with $q^2/N$ fixed.
\end{enumerate}


\section{Gaussian decay from time-independent infinite-temperature ensemble}

Consider a quantum system with Hamiltonian
\begin{equation}
    H(g)=\sum_I g_I V_I,
\end{equation}
where $I$ is a collective index, the $V_I$ are self-adjoint operators, and
the $g_I$ are real coupling parameters.  We choose the $g_I$ independently
from a Gaussian distribution,
\begin{equation}
    d\mu(g)
    =\prod_I
    \frac{dg_I}{\sqrt{2\pi\sigma}}
    e^{-g_I^2/(2\sigma)},
\end{equation}
so that
\begin{equation}
    \mathbb E_g[g_I]=0,
    \qquad
    \mathbb E_g[g_Ig_J]=\sigma\delta_{IJ}.
\end{equation}
For the time-independent problem $\sigma$ has units of energy squared.

We work throughout this section at infinite temperature.  Fix once and for
all an orthonormal basis $\{|i\rangle\}_{i=1}^d$ of the microscopic Hilbert
space and define
\begin{equation}
    |1\rangle
    :=\frac{1}{\sqrt d}\sum_{i=1}^d|i\rangle_L|i\rangle_R.
\end{equation}
For every operator $O$,
\begin{equation}
    (O_L-O_R^T)|1\rangle=0,
\end{equation}
where the transpose is taken in the fixed basis used to define $|1\rangle$.
For self-adjoint $O$, $O^T=O^*$.

For fixed $g$, the Heisenberg evolution is
\begin{equation}
    O(t)=e^{-iH(g)t}Oe^{iH(g)t}.
\end{equation}
On the doubled Hilbert space this becomes
\begin{equation}
    O(t)|1\rangle
    =e^{-it\widehat H(g)}O|1\rangle,
    \qquad
    \widehat H(g)=\sum_I g_I\widehat V_I,
    \qquad
    \widehat V_I:=V_{I,L}-V_{I,R}^T.
\end{equation}
The disorder-averaged evolution is the unital completely positive map
\begin{equation}
    \mathbb E_g[O(t)]
    =\int d\mu(g)\,
      e^{-iH(g)t}Oe^{iH(g)t},
\end{equation}
which on the doubled space is represented by
\begin{equation}
    \mathbb E_g[O(t)]|1\rangle
    =\int d\mu(g)\,
      e^{-it\sum_I g_I\widehat V_I}O|1\rangle.
\end{equation}
Finite-temperature correlators and their GNS representation are reviewed
separately in Appendix~A.  Their quenched disorder average involves the
random thermal normalization and will not be needed in this section.


\subsection{Gaussian average and ordered Wick expansion}

When there is only a single coupling $g$, the Gaussian integral can be
performed explicitly,
\begin{equation}
    \int\frac{dg}{\sqrt{2\pi\sigma}}
    e^{-g^2/(2\sigma)}e^{-itg\widehat V}
    =e^{-\frac{\sigma t^2}{2}\widehat V^2}.
\end{equation}
More generally, if all the operators $\widehat V_I$ mutually commute, the
same calculation can be applied in a common eigenbasis and yields
\begin{equation}
    \int d\mu(g)\,
    e^{-it\sum_Ig_I\widehat V_I}
    =
    e^{-\frac{\sigma t^2}{2}\sum_I\widehat V_I^2}.
\end{equation}
Thus mutually commuting interactions produce an exact Gaussian decay.
For noncommuting $\widehat V_I$, the Gaussian integral does not reduce to
the exponential of $\sum_I\widehat V_I^2$; the deviations are controlled
by operator ordering in the Wick expansion.

Expanding the exponential and using Wick's theorem gives the exact ordered
pairing series
\begin{equation}
\begin{split}
    \int d\mu(g)\,
    e^{-it\sum_Ig_I\widehat V_I}
    ={}&
    \sum_{n=0}^{\infty}
    \frac{(-1)^n(\sigma t^2)^n}{(2n)!}
    \sum_{P\in{\cal P}_2(2n)}
    \sum_{I_1,\ldots,I_n}
    \widehat V_{I_{\pi_P(1)}}\cdots
    \widehat V_{I_{\pi_P(2n)}}.
\end{split}
\end{equation}
Here ${\cal P}_2(2n)$ is the set of pairings of $\{1,\ldots,2n\}$ and
$\pi_P$ repeats each label $I_r$ in the two slots paired by $P$.  The
operator ordering is inherited from the original product.  In particular,
\begin{align}
    \int d\mu(g)\,
    e^{-it\sum_Ig_I\widehat V_I}
    ={}&1-\frac{\sigma t^2}{2}\sum_I\widehat V_I^2
    \nonumber\\
    &+\frac{\sigma^2t^4}{24}
    \sum_{I,J}
    \left(
       \widehat V_I^2\widehat V_J^2
       +\widehat V_I\widehat V_J\widehat V_I\widehat V_J
       +\widehat V_I\widehat V_J^2\widehat V_I
    \right)
    +O(t^6).
\end{align}
As a check, when all the $\widehat V_I$ commute the three fourth-order terms
coincide and reproduce the $1/8$ coefficient in the Taylor expansion of the
commuting Gaussian result.

It is useful to write
\begin{equation}
    \int d\mu(g)\,
    e^{-it\sum_Ig_I\widehat V_I}
    =e^{\Omega(t)},
    \qquad
    \Omega(t)=\sum_{n\ge1}\Omega_n(t),
\end{equation}
where $\Omega_n$ starts at order $t^{2n}$.  The terms in $\log M$ are
connected noncommutative combinations of the $\widehat V_I$.  The Gaussian
average can equivalently be written as the heat-kernel identity
\begin{equation}
    \int d\mu(g)\,f(g)
    =
    \left.
    \exp\!\left(
       \frac{\sigma}{2}
       \sum_I\frac{\partial^2}{\partial g_I^2}
    \right)f(g)
    \right|_{g=0}.
\end{equation}

\paragraph{The $k$-twirl channel and its moment operator.}
For a fixed realization $J$, define $U_J(t)=e^{-itH(J)}$ and
\begin{equation}
    {\cal U}_{J,t}^{(k)}(Y)
    :=U_J(t)^{\otimes k}Y U_J(t)^{\dagger\otimes k},
    \qquad
    Y\in{\rm End}({\cal H}^{\otimes k}).
\end{equation}
The disorder-averaged $k$-twirl map is
\begin{equation}
    {\cal E}_t^{(k)}(Y)
    :=\mathbb E_J\!\left[{\cal U}_{J,t}^{(k)}(Y)\right].
\end{equation}
Using vectorization $|Y\rangle\rangle=(Y\otimes1)|1\rangle$ on the doubled
$k$-replica space,
\begin{equation}
    (O_L-O_R^T)|Y\rangle\rangle
    =|[O,Y]\rangle\rangle.
\end{equation}
The corresponding moment operator is
\begin{equation}
    M_t^{(k)}
    :=\mathbb E_J\!\left[
       U_J(t)^{\otimes k}\otimes(U_J(t)^*)^{\otimes k}
    \right],
    \qquad
    |{\cal E}_t^{(k)}(Y)\rangle\rangle
    =M_t^{(k)}|Y\rangle\rangle.
\end{equation}
For each replica define
\begin{equation}
    \widehat V_I^{(a)}
    :=V_{I,L}^{(a)}-(V_{I,R}^{(a)})^T,
    \qquad
    \widehat V_I^{(k)}
    :=\sum_{a=1}^k\widehat V_I^{(a)}.
\end{equation}
Then
\begin{equation}
    M_t^{(k)}
    =\mathbb E_J\!\left[
       e^{-it\sum_IJ_I\widehat V_I^{(k)}}
    \right].
\end{equation}
Its ordered Wick series is therefore
\begin{equation}
    M_t^{(k)}
    =
    \sum_{n=0}^{\infty}
    \frac{(-1)^n(\sigma t^2)^n}{(2n)!}
    \sum_{P\in{\cal P}_2(2n)}
    \sum_{I_1,\ldots,I_n}
    \widehat V_{I_{\pi_P(1)}}^{(k)}\cdots
    \widehat V_{I_{\pi_P(2n)}}^{(k)}.
\end{equation}
In particular,
\begin{equation}
    M_t^{(k)}
    =1-\frac{\sigma t^2}{2}
    \sum_I(\widehat V_I^{(k)})^2+O(t^4).
\end{equation}


\subsection{Gaussian decay from large $N$ limit}

Now consider a system comprised of $N$ degrees of freedom,
$i=1,\ldots,N$, with $q$-body interaction terms.  The label $I$ is a set
of $q$ distinct sites,
\begin{equation}
    I=\{i_1,\ldots,i_q\},
\end{equation}
and we write
\begin{equation}
    H
    =\sum_{i_1<\cdots<i_q}
    g_{i_1\cdots i_q}V_{i_1\cdots i_q}.
\end{equation}
If the local operators on distinct sites commute, then interactions with
disjoint support commute exactly.

For comparison, if there is a Majorana fermion at each site,
\begin{equation}
    \{\chi_i,\chi_j\}=2\delta_{ij},
\end{equation}
and $V_I=\chi_{i_1}\cdots\chi_{i_q}$ with $q$ even, then two $q$-body
monomials satisfy
\begin{equation}
    [V_I,V_J]
    =\left(1-(-1)^{|I\cap J|}\right)V_IV_J.
\end{equation}
Thus they commute for an even number of overlaps and anticommute for an odd
number.  We return to fermionic SYK in Sec.~\ref{sec:largeN_SYK}.

Choosing two random $q$-element subsets $I,J$, the probability of at least
one overlap is
\begin{equation}
    P({\rm overlap})
    =1-\frac{\binom{N-q}{q}}{\binom Nq}
    =\frac{q^2}{N}+O(N^{-2})
\end{equation}
for fixed $q$ and large $N$.  Therefore almost all pairs commute in the
large-$N$ limit.  The vanishing overlap probability alone is not sufficient,
however, because the number of interaction terms grows with $N$.  We must
combine it with the connected cumulant expansion.


\subsection{Gaussian decay for the $k$-twirl at large $N$}
\label{subsec:quenched_syk_gaussian_regulator}

Consider
\begin{equation}
    H_N(J)=\sum_{I\in{\cal A}_N}J_IV_I,
    \qquad
    V_I=V_I^\dagger,
\end{equation}
with independent centered Gaussian couplings
\begin{equation}
    \mathbb E[J_I]=0,
    \qquad
    \mathbb E[J_IJ_J]=\sigma_N\delta_{IJ}.
\end{equation}
We keep $q$ and $k$ fixed as $N\to\infty$ and assume
\begin{equation}
    M_N:=|{\cal A}_N|=\Theta(N^q),
    \qquad
    \sigma_N=\Theta(N^{-(q-1)}),
\end{equation}
with noncommutation requiring overlaps of fixed-size supports.  Define
\begin{equation}
    \widehat V_I^{(k)}
    =\sum_{a=1}^k
    \left(V_{I,L}^{(a)}-(V_{I,R}^{(a)})^T\right).
\end{equation}
Then
\begin{equation}
    M_t^{(k)}
    =\mathbb E_J\!\left[
       e^{-it\sum_IJ_I\widehat V_I^{(k)}}
    \right]
    =e^{\Omega(t)}.
\end{equation}
The leading term is
\begin{equation}
    \Omega_1(t)
    =-\frac{\sigma_Nt^2}{2}
    \sum_I(\widehat V_I^{(k)})^2.
\end{equation}
At fourth order,
\begin{align}
    \Omega_2(t)
    =\frac{\sigma_N^2t^4}{24}
    \sum_{I,J}
    \bigg[&
       \widehat V_I^{(k)}\widehat V_J^{(k)}
       \widehat V_I^{(k)}\widehat V_J^{(k)}
       +\widehat V_I^{(k)}(\widehat V_J^{(k)})^2
        \widehat V_I^{(k)}
       \nonumber\\
       &-2(\widehat V_I^{(k)})^2(\widehat V_J^{(k)})^2
    \bigg].
\end{align}
The bracket vanishes whenever
$[\widehat V_I^{(k)},\widehat V_J^{(k)}]=0$.  Thus only an $O(1/N)$
fraction of the $M_N^2$ pairs contribute.

The natural diffusive variable is
\begin{equation}
    \tau_N:=\sigma_NM_Nt^2=O(1),
    \qquad
    t^2=O(N^{-1}).
\end{equation}
Then $\Omega_1=O(1)$ and
\begin{equation}
    \Omega_2
    =O\!\left(
       \sigma_N^2t^4M_N^2\frac1N
    \right)
    =O(N^{-1}).
\end{equation}

The same counting extends order by order in $\log M_t^{(k)}$.  Consider a
term in $\Omega_n$ containing $r\le n$ distinct interaction labels.  Join
two labels when the corresponding hatted operators fail to commute.  If
this noncommutation graph is disconnected, the Gaussian average factorizes
between commuting components and the disconnected contribution is removed
by the logarithm.  Hence a contribution to $\Omega_n$ has a connected
noncommutation graph.

For fixed $q$ and $r$, after choosing the first interaction freely, every
additional vertex in a spanning tree must overlap one of the previously
chosen supports, which costs a factor $O(1/N)$.  The number of connected
$r$-tuples is therefore
\begin{equation}
    O\!\left(\frac{M_N^r}{N^{r-1}}\right).
\end{equation}
Consequently,
\begin{align}
    \Omega_n
    &=O\!\left[
       (\sigma_Nt^2)^n
       \frac{M_N^r}{N^{r-1}}
    \right]
    \nonumber\\
    &=O(N^{-(n-1)}),
    \qquad n\ge2,
\end{align}
where we used $M_N=\Theta(N^q)$ and
$q(n-r)+(r-1)\ge n-1$.

Thus, order by order in the fixed-$q$, fixed-$k$ large-$N$ expansion,
\begin{equation}
    M_t^{(k)}
    \simeq
    \exp\!\left[
       -\frac{\sigma_Nt^2}{2}
       \sum_I(\widehat V_I^{(k)})^2
    \right].
\end{equation}
The counting establishes the suppression at every fixed order in $\tau_N$.
A uniform estimate for the fully resummed series at arbitrary finite
$\tau_N$ would require an additional summability bound.

Define
\begin{equation}
    \widehat X_I
    :=\sqrt{\sigma_N}\,\widehat V_I^{(k)},
    \qquad
    \widehat G_{\rm eff}^{(k)}
    :=\frac12\sum_I\widehat X_I^2.
\end{equation}
Then
\begin{equation}
    \boxed{
    M_t^{(k)}\simeq e^{-t^2\widehat G_{\rm eff}^{(k)}}.
    }
\end{equation}
On the physical $k$-replica space define
\begin{equation}
    X_I
    :=\sqrt{\sigma_N}\sum_{a=1}^kV_I^{(a)}.
\end{equation}
The corresponding superoperator is
\begin{align}
    G_{\rm eff}^{(k)}(Y)
    &=\frac12\sum_I[X_I,[X_I,Y]]
    \nonumber\\
    &=-\sum_I
    \left(
       X_IYX_I-\frac12\{X_I^2,Y\}
    \right),
\end{align}
so that
\begin{equation}
    \boxed{
    {\cal E}_t^{(k)}\simeq e^{-t^2G_{\rm eff}^{(k)}}.
    }
\end{equation}


\section{Gaussian decay in large $N$ bosonic and Spin SYK}
\label{sec:bosonic_spin_syk_gaussian_regulator}

In Sec.~\ref{subsec:quenched_syk_gaussian_regulator} we found that the
quenched $k$-twirl becomes a Gaussian decay at large $N$ when
noncommutation requires rare overlaps.  We now apply this result to the
bosonic Spin SYK models of \cite{swingle2024bosonic} and the two-component Spin
SYK models of \cite{basu2026complexity}.  Operators on different physical sites
commute exactly, so the overlap counting is particularly direct.

For a Hamiltonian
\begin{equation}
    H(J)=\sum_{I\in{\cal A}_N}J_IV_I,
    \qquad
    \mathbb E[J_IJ_J]=\sigma_N\delta_{IJ},
\end{equation}
define
\begin{equation}
    W_I^{(k)}:=\sum_{a=1}^kV_I^{(a)}.
\end{equation}
The leading generator is
\begin{equation}
    G_{\rm eff}^{(k)}(Y)
    =\frac{\sigma_N}{2}
    \sum_{I\in{\cal A}_N}
    [W_I^{(k)},[W_I^{(k)},Y]].
\end{equation}

\subsection{Three-component bosonic Spin SYK}

There are $N$ qubits and the interaction operators are
\begin{equation}
    V_I
    =\sigma_{r_1}^{\mu_1}\cdots\sigma_{r_q}^{\mu_q},
    \qquad
    r_1<\cdots<r_q,
    \qquad
    \mu_j\in\{x,y,z\}.
\end{equation}
The Hamiltonian is
\begin{equation}
    H
    =\sum_{\substack{r_1<\cdots<r_q\\
                      \mu_1,\ldots,\mu_q}}
    J_{r_1\mu_1\cdots r_q\mu_q}
    \sigma_{r_1}^{\mu_1}\cdots\sigma_{r_q}^{\mu_q},
\end{equation}
with
\begin{equation}
    \mathbb E[J_I]=0,
    \qquad
    \mathbb E[J_IJ_J]=\sigma_q^{(3)}\delta_{IJ},
    \qquad
    \sigma_q^{(3)}=\frac{(q-1)!J^2}{N^{q-1}}.
\end{equation}
There are
\begin{equation}
    M_q^{(3)}=3^q\binom Nq
\end{equation}
interaction terms.

Suppose two interaction terms overlap on $m$ sites.  At each shared site
the two Pauli matrices disagree with probability $2/3$.  If $d$ is the
number of shared sites with different Pauli components,
\begin{equation}
    V_IV_J=(-1)^dV_JV_I.
\end{equation}
Conditioned on $|I\cap J|=m$,
\begin{equation}
    P([V_I,V_J]\neq0\mid m)
    =\frac{1-(-1/3)^m}{2}.
\end{equation}
Hence
\begin{equation}
    p_{\rm nc}^{(3)}
    =\sum_{m=1}^q
    \frac{\binom qm\binom{N-q}{q-m}}{\binom Nq}
    \frac{1-(-1/3)^m}{2}
    =\frac{2q^2}{3N}+O(N^{-2}).
\end{equation}

The effective generator is
\begin{equation}
    G_{{\rm eff},q}^{(k,3)}(Y)
    =\frac{(q-1)!J^2}{2N^{q-1}}
    \sum_{\substack{r_1<\cdots<r_q\\
                    \mu_1,\ldots,\mu_q}}
    [W_I^{(k)},[W_I^{(k)},Y]],
\end{equation}
where
\begin{equation}
    W_I^{(k)}
    =\sum_{a=1}^k
    \left(
       \sigma_{r_1}^{\mu_1}\cdots\sigma_{r_q}^{\mu_q}
    \right)^{(a)}.
\end{equation}
The diffusive variable is
\begin{equation}
    \tau_q^{(3)}
    :=\sigma_q^{(3)}M_q^{(3)}t^2
    =\frac{3^q}{q}J^2Nt^2[1+O(N^{-1})].
\end{equation}
For $q=2$,
\begin{equation}
    \sigma_2^{(3)}=\frac{J^2}{N},
    \qquad
    M_2^{(3)}=9\binom N2,
    \qquad
    p_{\rm nc}^{(3)}=\frac{8}{3N}+O(N^{-2}),
\end{equation}
and
\begin{equation}
    G_{{\rm eff},2}^{(k,3)}(Y)
    =\frac{J^2}{2N}
    \sum_{r<s}\sum_{\mu,\nu=x,y,z}
    [W_{rs}^{\mu\nu},[W_{rs}^{\mu\nu},Y]],
\end{equation}
with
\begin{equation}
    W_{rs}^{\mu\nu}
    =\sum_{a=1}^k(\sigma_r^\mu\sigma_s^\nu)^{(a)}.
\end{equation}
For $q=4$,
\begin{equation}
    \sigma_4^{(3)}=\frac{6J^2}{N^3},
    \qquad
    M_4^{(3)}=81\binom N4,
    \qquad
    p_{\rm nc}^{(3)}=\frac{32}{3N}+O(N^{-2}),
\end{equation}
and
\begin{equation}
    G_{{\rm eff},4}^{(k,3)}(Y)
    =\frac{3J^2}{N^3}
    \sum_{r_1<r_2<r_3<r_4}
    \sum_{\mu_1,\ldots,\mu_4=x,y,z}
    [W_I^{(k)},[W_I^{(k)},Y]].
\end{equation}
Thus, for $q=2,4$,
\begin{equation}
    {\cal E}_{t,q}^{(k)}
    \simeq e^{-t^2G_{{\rm eff},q}^{(k,3)}},
    \qquad t^2=O(N^{-1}).
\end{equation}
This is a statement about the real-time quenched twirl in the diffusive
large-$N$ limit and is distinct from the low-temperature saddle-point
question studied in \cite{swingle2024bosonic}.

\subsection{Two-component Spin SYK}

Let $N_{\rm Spin}$ be the number of physical spin sites.  The
$2N_{\rm Spin}$ constituent operators are
\begin{equation}
    O_{2r-1}=\sigma_r^x,
    \qquad
    O_{2r}=\sigma_r^y,
    \qquad
    r=1,\ldots,N_{\rm Spin}.
\end{equation}
The Hamiltonian of \cite{basu2026complexity} is
\begin{equation}
    H_{{\rm Spin\,SYK}_q}
    =\sqrt{\frac{(q-1)!}{(2N_{\rm Spin})^{q-1}}}
    \sum_{1\le i_1<\cdots<i_q\le2N_{\rm Spin}}
    J_{i_1\cdots i_q}V_{i_1\cdots i_q},
\end{equation}
where
\begin{equation}
    V_{i_1\cdots i_q}
    :=i^{\eta_{i_1\cdots i_q}}O_{i_1}\cdots O_{i_q}
    =V_{i_1\cdots i_q}^\dagger
\end{equation}
and the $J_I$ are independent unit-variance Gaussians.  Equivalently,
\begin{equation}
    \sigma_q^{(2)}
    =\frac{(q-1)!}{(2N_{\rm Spin})^{q-1}}.
\end{equation}

For the genuine model, gSpin-SYK, no interaction contains two constituent
operators from the same physical site.  A genuine interaction is
\begin{equation}
    V_{S,\boldsymbol\alpha}
    =\prod_{r\in S}\sigma_r^{\alpha_r},
    \qquad
    |S|=q,
    \qquad
    \alpha_r\in\{x,y\},
\end{equation}
with
\begin{equation}
    M_{q,g}^{(2)}=2^q\binom{N_{\rm Spin}}q.
\end{equation}
If two genuine interactions overlap on $m\ge1$ physical sites, the number
of Pauli mismatches has odd parity with probability $1/2$.  Therefore
\begin{equation}
    p_{{\rm nc},g}^{(2)}
    =\frac12\left[
       1-\frac{\binom{N_{\rm Spin}-q}{q}}
               {\binom{N_{\rm Spin}}q}
    \right]
    =\frac{q^2}{2N_{\rm Spin}}+O(N_{\rm Spin}^{-2}).
\end{equation}
The effective generator is
\begin{equation}
    G_{{\rm eff},q,g}^{(k,2)}(Y)
    =\frac{(q-1)!}{2(2N_{\rm Spin})^{q-1}}
    \sum_{\substack{|S|=q\\
                    \boldsymbol\alpha\in\{x,y\}^q}}
    [W_{S,\boldsymbol\alpha}^{(k)},
     [W_{S,\boldsymbol\alpha}^{(k)},Y]],
\end{equation}
where
\begin{equation}
    W_{S,\boldsymbol\alpha}^{(k)}
    =\sum_{a=1}^k
    \left(\prod_{r\in S}\sigma_r^{\alpha_r}\right)^{(a)}.
\end{equation}
Its diffusive variable is
\begin{equation}
    \tau_{q,g}^{(2)}
    =\sigma_q^{(2)}M_{q,g}^{(2)}t^2
    =\frac{2}{q}N_{\rm Spin}t^2
      [1+O(N_{\rm Spin}^{-1})].
\end{equation}
For $q=2$,
\begin{equation}
    \sigma_2^{(2)}=\frac{1}{2N_{\rm Spin}},
    \qquad
    M_{2,g}^{(2)}=4\binom{N_{\rm Spin}}2,
    \qquad
    p_{{\rm nc},g}^{(2)}
    =\frac{2}{N_{\rm Spin}}+O(N_{\rm Spin}^{-2}),
\end{equation}
and
\begin{equation}
    G_{{\rm eff},2,g}^{(k,2)}(Y)
    =\frac{1}{4N_{\rm Spin}}
    \sum_{r<s}\sum_{\alpha,\beta=x,y}
    [W_{rs}^{\alpha\beta},[W_{rs}^{\alpha\beta},Y]].
\end{equation}
For $q=4$,
\begin{equation}
    \sigma_4^{(2)}=\frac{3}{4N_{\rm Spin}^3},
    \qquad
    M_{4,g}^{(2)}=16\binom{N_{\rm Spin}}4,
    \qquad
    p_{{\rm nc},g}^{(2)}
    =\frac{8}{N_{\rm Spin}}+O(N_{\rm Spin}^{-2}),
\end{equation}
and
\begin{equation}
    G_{{\rm eff},4,g}^{(k,2)}(Y)
    =\frac{3}{8N_{\rm Spin}^3}
    \sum_{r_1<r_2<r_3<r_4}
    \sum_{\alpha_1,\ldots,\alpha_4=x,y}
    [W_I^{(k)},[W_I^{(k)},Y]].
\end{equation}
Thus
\begin{equation}
    {\cal E}_{t,q,g}^{(k)}
    \simeq e^{-t^2G_{{\rm eff},q,g}^{(k,2)}},
    \qquad t^2=O(N_{\rm Spin}^{-1}).
\end{equation}

\subsection{Full versus genuine Spin SYK}

The full Spin-SYK model allows both $O_{2r-1}$ and $O_{2r}$ from the same
physical site.  The total number of interaction terms is
\begin{equation}
    M_{q,{\rm full}}^{(2)}=\binom{2N_{\rm Spin}}q,
\end{equation}
whereas the number of genuine terms is
$2^q\binom{N_{\rm Spin}}q$.  Their ratio is
\begin{equation}
    \frac{M_{q,g}^{(2)}}{M_{q,{\rm full}}^{(2)}}
    =1-\frac{q(q-1)}{4N_{\rm Spin}}
     +O(N_{\rm Spin}^{-2}).
\end{equation}
Thus the number of collision terms is $O(N_{\rm Spin}^{q-1})$.  Since
$\sigma_q^{(2)}=O(N_{\rm Spin}^{-(q-1)})$, their contribution to the
unscaled generator is $O(1)$, whereas the genuine $q$-site sector is
$O(N_{\rm Spin})$.  On the diffusive scale,
\begin{equation}
    t^2\Delta G_{\rm eff}^{(k)}=O(N_{\rm Spin}^{-1}).
\end{equation}
Consequently,
\begin{equation}
    {\cal E}_{t,q,{\rm full}}^{(k)}
    \simeq {\cal E}_{t,q,g}^{(k)}
    \simeq e^{-t^2G_{{\rm eff},q,g}^{(k,2)}}.
\end{equation}
The self-site terms of the quadratic model and the lower-support collision
terms of the quartic model first affect the diffusive evolution at order
$1/N_{\rm Spin}$.

\subsection{Conserved $\mathbb Z_2$ charge in the two-component model}

Define
\begin{equation}
    \Gamma^3:=\bigotimes_{r=1}^{N_{\rm Spin}}\sigma_r^z.
\end{equation}
Every constituent operator anticommutes with $\Gamma^3$,
\begin{equation}
    \{\Gamma^3,O_i\}=0,
\end{equation}
so an interaction containing $q$ constituents satisfies
\begin{equation}
    \Gamma^3V_I=(-1)^qV_I\Gamma^3.
\end{equation}
For $q=2,4$,
\begin{equation}
    [\Gamma^3,V_I]=0.
\end{equation}
Hence each replica charge $\Gamma^{3(a)}$ commutes with every jump operator
\begin{equation}
    X_I=\sqrt{\sigma_q^{(2)}}\sum_{b=1}^kV_I^{(b)}
\end{equation}
and therefore lies in the kernel algebra of $G_{\rm eff}^{(k)}$.  The
two-component twirl cannot converge to the unrestricted Haar twirl on the
full Hilbert space.  Determining the complete invariant algebra requires a
separate analysis of the interaction algebra.

\subsection{Summary}

The quadratic and quartic bosonic Spin SYK models obey the same large-$N$
mechanism as Sec.~\ref{subsec:quenched_syk_gaussian_regulator}.  In each
case $M_N=\Theta(N^q)$, $\sigma_NM_N=\Theta(N)$, and the probability of a
nonzero commutator is $O(1/N)$.  Thus
\begin{equation}
    \tau_N=\sigma_NM_Nt^2=O(1),
    \qquad
    t=O(N^{-1/2}),
\end{equation}
and
\begin{equation}
    \boxed{
    {\cal E}_t^{(k)}\simeq e^{-t^2G_{\rm eff}^{(k)}},
    \qquad
    G_{\rm eff}^{(k)}(Y)
    =\frac{\sigma_N}{2}\sum_I
    \left[
       \sum_{a=1}^kV_I^{(a)},
       \left[\sum_{b=1}^kV_I^{(b)},Y\right]
    \right].
    }
\end{equation}


\section{Gaussian decay in large $N$ SYK}
\label{sec:largeN_SYK}

We now return to fermionic SYK and study the structure of the effective
generator in more detail.  The model contains $N$ Majorana fermions,
\begin{equation}
    \{\chi_i,\chi_j\}=2\delta_{ij},
\end{equation}
with Hamiltonian
\begin{equation}
    H
    =i^{q/2}\sum_{i_1<\cdots<i_q}
    J_{i_1\cdots i_q}\chi_{i_1}\cdots\chi_{i_q}.
\end{equation}
We restrict to even $q$ and even $N$.  Define
\begin{equation}
    V_I:=\chi_{i_1}\cdots\chi_{i_q},
    \qquad
    W_I:=i^{q/2}V_I.
\end{equation}
Then
\begin{equation}
    V_I^2=(-1)^{q(q-1)/2}=(-1)^{q/2},
    \qquad
    W_I=W_I^\dagger,
    \qquad
    W_I^2=1.
\end{equation}
The static Gaussian couplings satisfy
\begin{equation}
    \mathbb E[J_I]=0,
    \qquad
    \mathbb E[J_IJ_J]=\sigma\delta_{IJ},
    \qquad
    \sigma=\frac{2^{q-1}q!}{q^2N^{q-1}}J.
\end{equation}
For $q=2$ and $q=4$ this gives
\begin{equation}
    \sigma=\frac{J}{N},
    \qquad
    \sigma=\frac{12J}{N^3},
\end{equation}
respectively.

\subsection{Jordan-Wigner transform}

Pair $\chi_{2j-1},\chi_{2j}$ into complex fermions
\begin{equation}
    c_j:=\frac12(\chi_{2j-1}+i\chi_{2j}),
\end{equation}
with
\begin{equation}
    \{c_l,c_j^\dagger\}=\delta_{lj},
    \qquad
    \{c_l,c_j\}=\{c_l^\dagger,c_j^\dagger\}=0.
\end{equation}
Use the standard qubit convention
\begin{equation}
    \sigma_j^-:=\frac12(X_j+iY_j)=|0\rangle\langle1|_j,
    \qquad
    N_j:=c_j^\dagger c_j=\frac{1-Z_j}{2}.
\end{equation}
Then
\begin{equation}
    Z_j=-i\chi_{2j-1}\chi_{2j},
    \qquad
    (-1)^{N_j}=Z_j,
\end{equation}
and
\begin{align}
    c_j&=(Z_1\cdots Z_{j-1})\sigma_j^-,
    \nonumber\\
    \chi_{2j-1}&=(Z_1\cdots Z_{j-1})X_j,
    \nonumber\\
    \chi_{2j}&=(Z_1\cdots Z_{j-1})Y_j.
\end{align}
Let
\begin{equation}
    |e\rangle
    =\frac1{\sqrt d}\sum_{\mathbf n}
    |\mathbf n\rangle_L|\mathbf n\rangle_R
\end{equation}
be the maximally entangled vector in this fixed Jordan-Wigner basis.  For
any operator $O$,
\begin{equation}
    (O_L-O_R^T)|e\rangle=0.
\end{equation}
Define
\begin{equation}
    \chi_{j,R}:=(\chi_j^T)_R.
\end{equation}
Then
\begin{equation}
    (\chi_{j,L}-\chi_{j,R})|e\rangle=0,
\end{equation}
while $L$ and $R$ Majoranas commute because they act on different tensor
factors.  For the Hermitian interaction $W_I$,
\begin{align}
    (W_I^T)_R
    &=i^{q/2}\chi_{i_q,R}\cdots\chi_{i_1,R}
    \nonumber\\
    &=(-i)^{q/2}\chi_{i_1,R}\cdots\chi_{i_q,R}.
\end{align}

\subsection{Totally anti-commuting variables}

For the $k$-twirl we have $k$ $L$ copies and $k$ $R$ copies,
\begin{equation}
    \chi_{l,L}^{(a)},\qquad
    \chi_{l,R}^{(a)},
    \qquad l=1,\ldots,N,\quad a=1,\ldots,k.
\end{equation}
Majoranas on different tensor factors commute.  Introduce the parity on
each copy,
\begin{equation}
    Q_{L/R}^{(a)}
    :=(-1)^{\sum_{j=1}^{N/2}N_j^{(a)}}
    =\prod_{j=1}^{N/2}Z_{j,L/R}^{(a)}
    =(-i)^{N/2}\prod_{l=1}^N\chi_{l,L/R}^{(a)}.
\end{equation}
It obeys
\begin{equation}
    \{Q_{L/R}^{(a)},\chi_{l,L/R}^{(a)}\}=0,
    \qquad
    (Q_{L/R}^{(a)})^2=1.
\end{equation}
Define
\begin{equation}
    \Sigma^{(a)}
    :=\left(\prod_{b=1}^{a-1}Q_L^{(b)}Q_R^{(b)}\right)Q_L^{(a)},
\end{equation}
and
\begin{equation}
    \psi_{l,L}^{(a)}:=i\Sigma^{(a)}\chi_{l,L}^{(a)},
    \qquad
    \psi_{l,R}^{(a)}:=\Sigma^{(a)}\chi_{l,R}^{(a)}.
\end{equation}
These are self-adjoint Majoranas satisfying
\begin{equation}
    \{\psi_{l,\alpha}^{(a)},\psi_{l',\beta}^{(b)}\}
    =2\delta_{ll'}\delta_{ab}\delta_{\alpha\beta}.
\end{equation}
For even $q$ the Klein factors cancel inside a $q$-body monomial,
\begin{align}
    \chi_{i_1,L}^{(a)}\cdots\chi_{i_q,L}^{(a)}
    &=\psi_{i_1,L}^{(a)}\cdots\psi_{i_q,L}^{(a)},
    \nonumber\\
    \chi_{i_1,R}^{(a)}\cdots\chi_{i_q,R}^{(a)}
    &=\psi_{i_1,R}^{(a)}\cdots\psi_{i_q,R}^{(a)}.
\end{align}
Define
\begin{equation}
    \Psi_{I,\alpha}^{(a)}
    :=\psi_{i_1,\alpha}^{(a)}\cdots\psi_{i_q,\alpha}^{(a)}.
\end{equation}
The hatted interaction is
\begin{equation}
    \widehat W_I^{(k)}
    =i^{q/2}\sum_{a=1}^k\Psi_{I,L}^{(a)}
     -(-i)^{q/2}\sum_{a=1}^k\Psi_{I,R}^{(a)}.
\end{equation}
Bundle $(a,\alpha)$ into a flavor index $A=1,\ldots,2k$ and define
\begin{equation}
    s_A=
    \begin{cases}
       1,&\alpha=L,\\
       -(-1)^{q/2},&\alpha=R.
    \end{cases}
\end{equation}
Then
\begin{equation}
    \widehat W_I^{(k)}
    =i^{q/2}\sum_{A=1}^{2k}s_A\Psi_I^A.
\end{equation}

\subsection{$k$-twirl map in large $N$ SYK}

From Sec.~\ref{subsec:quenched_syk_gaussian_regulator},
\begin{equation}
    \widehat G_{\rm eff}^{(k)}
    =\frac{\sigma}{2}
    \sum_{i_1<\cdots<i_q}
    (\widehat W_{i_1\cdots i_q}^{(k)})^2.
\end{equation}
Define antisymmetric bilinears
\begin{equation}
    M_i^{AB}:=\frac{i}{2}[\psi_i^A,\psi_i^B],
    \qquad
    M_i^{AB}=-M_i^{BA}.
\end{equation}
For $A\neq B$,
\begin{equation}
    M_i^{AB}=i\psi_i^A\psi_i^B,
    \qquad
    (M_i^{AB})^\dagger=M_i^{AB},
    \qquad
    (M_i^{AB})^2=1,
\end{equation}
while $M_i^{AA}=0$.  For $i\neq j$,
\begin{equation}
    [M_i^{AB},M_j^{CD}]=0,
\end{equation}
and at the same site
\begin{equation}
    [M_i^{AB},\psi_j^C]
    =2i\delta_{ij}
    (\delta^{BC}\psi_i^A-\delta^{AC}\psi_i^B),
\end{equation}
\begin{equation}
    [M_i^{AB},M_j^{CD}]
    =2i\delta_{ij}
    (\delta^{BC}M_i^{AD}-\delta^{AC}M_i^{BD}
     -\delta^{BD}M_i^{AC}+\delta^{AD}M_i^{BC}).
\end{equation}
Thus every site carries a copy of $\mathfrak{so}(2k)$.

For $A\neq B$ and even $q$,
\begin{equation}
    \Psi_I^A\Psi_I^B
    =M_{i_1}^{AB}\cdots M_{i_q}^{AB},
    \qquad
    i^q(\Psi_I^A)^2=1.
\end{equation}
Therefore
\begin{equation}
    \boxed{
    \widehat G_{\rm eff}^{(k)}
    =\sigma k\binom Nq
     +\sigma i^q\sum_{A<B}s_As_B e_q^{AB},
    }
\end{equation}
where
\begin{equation}
    e_q^{AB}
    :=\sum_{i_1<\cdots<i_q}
    M_{i_1}^{AB}\cdots M_{i_q}^{AB}.
\end{equation}

\subsection{Properties of the generator $G_{\rm eff}^{(k)}$}

Fix $A<B$ and define
\begin{equation}
    S^{AB}:=\sum_{i=1}^N M_i^{AB}.
\end{equation}
The $M_i^{AB}$ commute for different $i$ and square to one, so their joint
eigenvalues are $\pm1$ and
\begin{equation}
    S^{AB}\in\{-N,-N+2,\ldots,N-2,N\}.
\end{equation}
The number of sign configurations with a given value of $S^{AB}$ is
binomial; the full Hilbert-space multiplicity can contain an additional
flavor degeneracy that will not be important below.

The elementary symmetric polynomial $e_q^{AB}$ depends only on $S^{AB}$.
Its generating function is
\begin{equation}
    \sum_{q=0}^N z^q e_q^{AB}
    =\prod_{i=1}^N(1+zM_i^{AB})
    =(1+z)^{(N+S^{AB})/2}
     (1-z)^{(N-S^{AB})/2}.
\end{equation}
For even $q$,
\begin{equation}
    e_q^{AB}
    =\sum_{r=0}^{q/2}c_{q,r}(N)(S^{AB})^{q-2r}.
\end{equation}
For example,
\begin{equation}
    e_2^{AB}=\frac{(S^{AB})^2-N}{2},
\end{equation}
\begin{equation}
    e_4^{AB}
    =\frac{
       (S^{AB})^4-2(3N-4)(S^{AB})^2+3N(N-2)
    }{24}.
\end{equation}

Define the monic Krawtchouk normalization
\begin{equation}
    P_q(S):=q!\,e_q(S).
\end{equation}
Newton identities give
\begin{equation}
    P_{q+1}(S)
    =SP_q(S)-q(N-q+1)P_{q-1}(S),
    \qquad
    P_0(S)=1,\quad P_1(S)=S.
\end{equation}
Equivalently,
\begin{equation}
    P_q(S)=\det(S\mathbf 1_q-A_q),
\end{equation}
where $A_q$ is the $q\times q$ tridiagonal matrix with zero diagonal and
off-diagonal entries
\begin{equation}
    (A_q)_{m,m+1}=(A_q)_{m+1,m}
    =\sqrt{m(N-m+1)},
    \qquad m=1,\ldots,q-1.
\end{equation}
Expanding
\begin{equation}
    P_q(S)
    =\sum_{r=0}^{\lfloor q/2\rfloor}
    \widetilde c_{q,r}(N)S^{q-2r},
\end{equation}
we have $c_{q,r}=\widetilde c_{q,r}/q!$ and
\begin{equation}
    (q+1)c_{q+1,r}(N)
    =c_{q,r}(N)-(N-q+1)c_{q-1,r-1}(N).
\end{equation}
A closed form is
\begin{equation}
    c_{q,r}(N)
    =\frac{(-1)^r}{q!}
    \sum_{\substack{1\le m_1<\cdots<m_r\le q-1\\
                    m_{j+1}\ge m_j+2}}
    \prod_{j=1}^r m_j(N-m_j+1).
\end{equation}
In particular,
\begin{equation}
    c_{q,0}(N)=\frac1{q!},
    \qquad
    c_{q,1}(N)
    =-\frac{q(q-1)(3N-2q+4)}{6q!},
\end{equation}
and for $q=2n$ and even $N$,
\begin{equation}
    c_{2n,n}(N)=(-1)^n\binom{N/2}{n}.
\end{equation}

\subsubsection{Large-$N$ limits of $e_q^{AB}$}

Let
\begin{equation}
    m^{AB}:=\frac{S^{AB}}{N}.
\end{equation}
For $N\to\infty$ with $q$ fixed,
\begin{equation}
    \frac{e_q^{AB}}{\binom Nq}
    =(m^{AB})^q
     -\frac{q(q-1)}{2N}
      (1-(m^{AB})^2)(m^{AB})^{q-2}
     +O(N^{-2}).
\end{equation}
In particular,
\begin{equation}
    e_q^{AB}
    =\binom Nq(m^{AB})^q+O(N^{q-1}).
\end{equation}

There is also a standard Hermite limit when the degree grows slowly enough
that $q^2/N\to0$.  Expanding the generating function at small $z$ gives
\begin{equation}
    (1+z)^{(N+S^{AB})/2}
    (1-z)^{(N-S^{AB})/2}
    =\exp\!\left[
       S^{AB}z-\frac N2z^2
       +O(S^{AB}z^3)+O(Nz^4)
    \right].
\end{equation}
In the regime in which the higher terms are negligible this gives
\begin{equation}
    e_q^{AB}
    \sim\frac{N^{q/2}}{q!}
    {\rm He}_q\!\left(\frac{S^{AB}}{\sqrt N}\right),
\end{equation}
where
\begin{equation}
    e^{xt-t^2/2}
    =\sum_{q=0}^{\infty}{\rm He}_q(x)\frac{t^q}{q!}.
\end{equation}
The genuinely double-scaled regime $q^2/N=O(1)$ is more subtle: the
higher terms in the logarithm of the generating function need not be
negligible, so the ordinary Hermite formula above should not be used as a
uniform asymptotic statement in that regime.

Finally,
\begin{equation}
    \widehat G_{\rm eff}^{(k)}
    =\sigma k\binom Nq
     +\sigma i^q\sum_{A<B}s_As_B
      \sum_{r=0}^{q/2}c_{q,r}(N)(S^{AB})^{q-2r}.
\end{equation}
The global sums $S^{AB}=\sum_iM_i^{AB}$ generate the diagonal
$\mathfrak{so}(2k)$ action, and the doubled-space generator is an element
of $U(\mathfrak{so}(2k))$.

On the physical $k$-replica Hilbert space,
\begin{equation}
    X_I:=\sqrt\sigma\sum_{a=1}^kW_I^{(a)},
\end{equation}
and
\begin{equation}
    G_{\rm eff}^{(k)}(Y)
    =\frac12\sum_I[X_I,[X_I,Y]]
    =-\sum_I\left(X_IYX_I-\frac12\{X_I^2,Y\}\right).
\end{equation}
Thus $-G_{\rm eff}^{(k)}$ is in Lindblad form and
\begin{equation}
    {\cal E}_t^{(k)}\simeq e^{-t^2G_{\rm eff}^{(k)}}.
\end{equation}


\section{Exponential decay and semigroups from Brownian dynamics}

We now let the couplings fluctuate independently in time.  Divide a total
time interval $T$ into steps of width $\epsilon$ and take
\begin{equation}
    H_g(t_i)=\sum_I g_I(t_i)V_I,
\end{equation}
with independent Gaussian samples
\begin{equation}
    \mathbb E_g[g_I(t_i)]=0,
    \qquad
    \mathbb E_g[g_I(t_i)g_J(t_j)]
    =\frac{\sigma}{\epsilon}\,
     \delta_{IJ}\delta_{ij}.
\end{equation}
Here $\sigma$ is the Brownian noise strength.  Its dimensions differ from
the static variance used in Sec.~2.

For one time step,
\begin{equation}
    dU=U(t_i,t_i+\epsilon),
\end{equation}
and on the doubled $k$-replica space
\begin{equation}
    \widehat H_g^{(k)}
    :=\sum_{a=1}^k
    \left(H_{g,L}^{(a)}-(H_{g,R}^{(a)})^T\right)
    =\sum_Ig_I(t_i)\widehat V_I^{(k)}.
\end{equation}
Using a fixed vectorization basis, the infinitesimal moment is
\begin{align}
    \mathbb E_g\!\left[dU^{\otimes k}\otimes(dU^*)^{\otimes k}\right]
    &={\bf 1}
      -\frac{\epsilon^2}{2}
       \mathbb E_g[(\widehat H_g^{(k)})^2]
      +O(\epsilon^2)
    \nonumber\\
    &={\bf 1}-\epsilon\widehat G_{\rm eff}^{(k)}+O(\epsilon^2),
\end{align}
where odd Gaussian moments vanish and
\begin{equation}
    \boxed{
    \widehat G_{\rm eff}^{(k)}
    :=\frac{\epsilon}{2}
      \mathbb E_g[(\widehat H_g^{(k)})^2]
    =\frac{\sigma}{2}\sum_I(\widehat V_I^{(k)})^2.
    }
\end{equation}
Equivalently, define the self-adjoint operators
\begin{equation}
    X_I:=\sqrt\sigma\sum_{a=1}^kV_I^{(a)}.
\end{equation}
Then
\begin{equation}
    G_{\rm eff}^{(k)}(Y)
    =\frac12\sum_I[X_I,[X_I,Y]]
    =-\sum_I
      \left(X_IYX_I-\frac12\{X_I^2,Y\}\right).
\end{equation}
Thus the averaged evolution satisfies
\begin{equation}
    \partial_T\mathbb E_g[Y(T)]
    =-G_{\rm eff}^{(k)}\mathbb E_g[Y(T)].
\end{equation}
Taking $T/\epsilon$ independent time steps and sending $\epsilon\to0$ gives
the exact Markov semigroup
\begin{equation}
    \boxed{
    {\cal E}_{g;T}^{(k)}=e^{-TG_{\rm eff}^{(k)}}.
    }
\end{equation}
Unlike the quenched result of Sec.~2.3, no large-$N$ approximation is
needed: the temporal white noise makes the quadratic generator exact at
finite $N$.


\section{Semigroups in Brownian bosonic and Spin SYK models}
\label{sec:brownian_spin_syk}

We now apply the general Brownian result to the spin models of Sec.~3.
For later reference, define
\begin{equation}
    W_I^{(k)}:=\sum_{a=1}^kV_I^{(a)},
    \qquad
    \widehat W_I^{(k)}
    :=\sum_{a=1}^k
      \left(V_{I,L}^{(a)}-(V_{I,R}^{(a)})^T\right).
\end{equation}
Then
\begin{equation}
    \widehat G_{\rm eff}^{(k)}
    =\frac{\sigma}{2}\sum_I(\widehat W_I^{(k)})^2,
\end{equation}
and
\begin{equation}
    G_{\rm eff}^{(k)}(Y)
    =\frac{\sigma}{2}\sum_I
      [W_I^{(k)},[W_I^{(k)},Y]].
\end{equation}
The static and Brownian problems therefore have the same quadratic generator
form; the former gives $e^{-t^2G}$ asymptotically at large $N$, while the
latter gives $e^{-TG}$ exactly.

\subsection{Brownian three-component Spin SYK}

The Brownianized Swingle-Winer model is
\begin{equation}
    H_{\rm SW}(t)
    =\sum_{\substack{r_1<\cdots<r_q\\
                    \mu_1,\ldots,\mu_q}}
    J_{r_1\mu_1\cdots r_q\mu_q}(t)
    \sigma_{r_1}^{\mu_1}\cdots\sigma_{r_q}^{\mu_q},
\end{equation}
with
\begin{equation}
    \mathbb E[J_I(t)J_J(t')]
    =\sigma_q^{(3)}\delta_{IJ}\delta(t-t').
\end{equation}
Keeping the same $N$-scaling as the static model, write
\begin{equation}
    \sigma_q^{(3)}
    =\frac{(q-1)!{\cal J}_3}{N^{q-1}}.
\end{equation}
The parameter ${\cal J}_3$ sets the Brownian clock and should not be
identified dimensionally with the static $J^2$ without specifying units.
The exact generator is
\begin{equation}
    G_{{\rm SW},q}^{(k)}(Y)
    =\frac{\sigma_q^{(3)}}{2}
    \sum_{\substack{r_1<\cdots<r_q\\
                    \mu_1,\ldots,\mu_q}}
    [W_I^{(k)},[W_I^{(k)},Y]].
\end{equation}
For $q=2$,
\begin{equation}
    G_{{\rm SW},2}^{(k)}(Y)
    =\frac{{\cal J}_3}{2N}
    \sum_{r<s}\sum_{\mu,\nu=x,y,z}
    [W_{rs}^{\mu\nu},[W_{rs}^{\mu\nu},Y]],
\end{equation}
and for $q=4$,
\begin{equation}
    G_{{\rm SW},4}^{(k)}(Y)
    =\frac{3{\cal J}_3}{N^3}
    \sum_{r_1<r_2<r_3<r_4}
    \sum_{\mu_1,\ldots,\mu_4=x,y,z}
    [W_I^{(k)},[W_I^{(k)},Y]].
\end{equation}
In both cases
\begin{equation}
    {\cal E}_{{\rm SW};T}^{(k)}
    =e^{-TG_{{\rm SW},q}^{(k)}}.
\end{equation}

\subsubsection{The four-replica sector and operator-size dynamics}

For $k=2$, the doubled Hilbert space contains four replicas.  Reordering
them as forward, backward, forward, backward gives
\begin{equation}
    {\cal H}_4
    =\frac{\sigma_q^{(3)}}{2}
    \sum_A
    \left(
       V_A^a-V_A^{*,b}+V_A^c-V_A^{*,d}
    \right)^2.
\end{equation}
This is the four-replica Brownian Hamiltonian used in
\cite{XuBrownianSpin}.  Let $S$ be a Pauli string.  Since $V_A^2=1$, the
single-interaction contribution vanishes when $[V_A,S]=0$, while if
$\{V_A,S\}=0$,
\begin{equation}
    \frac12
    \left(
       V_A^a-V_A^{*,b}+V_A^c-V_A^{*,d}
    \right)^2
    |S\otimes S^\dagger\rangle
    =4|S\otimes S^\dagger\rangle
     -4|V_AS\otimes S^\dagger V_A\rangle.
\end{equation}
Thus every interaction that anticommutes with $S$ generates a classical
transition
\begin{equation}
    S\longrightarrow V_AS
\end{equation}
at rate
\begin{equation}
    4\sigma_q^{(3)}.
\end{equation}

Because the three-component ensemble contains all Pauli strings of size
$q$, the evolution closes on the ordinary operator size
\begin{equation}
    m=|{\rm supp}(S)|.
\end{equation}
Take $S=\sigma_1^z\cdots\sigma_m^z$ and parameterize the size change by
\begin{equation}
    \Delta m=q-2\ell+1,
    \qquad \ell=1,\ldots,q.
\end{equation}
The exact transition rate is
\begin{align}
    \gamma_\ell^{(3)}(m)
    =4\sigma_q^{(3)}\sum_{n_z}
    &2^{2\ell-2n_z-1}
     3^{n_z+q-2\ell+1}
     \binom{m}{n_z}
     \binom{N-m}{n_z+q-2\ell+1}
    \nonumber\\
    &\times
     \binom{m-n_z}{2\ell-2n_z-1},
\end{align}
where inadmissible binomial coefficients are defined to vanish.  Therefore
\begin{equation}
    \partial_TP(m,T)
    =\sum_{\ell=1}^q
    \left[
       \gamma_\ell^{(3)}(m+2\ell-1-q)
       P(m+2\ell-1-q,T)
       -\gamma_\ell^{(3)}(m)P(m,T)
    \right].
\end{equation}
This is the spin master equation of \cite{XuBrownianSpin}.

For $q=2$, $\Delta m=\pm1$ and
\begin{equation}
    \gamma_1^{(3)}(m)
    =24\sigma_2^{(3)}m(N-m),
    \qquad
    \gamma_2^{(3)}(m)
    =8\sigma_2^{(3)}m(m-1).
\end{equation}
For $q=4$, $\Delta m=3,1,-1,-3$ and
\begin{align}
    \gamma_1^{(3)}(m)
    &=216\sigma_4^{(3)}m\binom{N-m}{3},
    \\
    \gamma_2^{(3)}(m)
    &=96\sigma_4^{(3)}\binom m3(N-m)
      +72\sigma_4^{(3)}m(m-1)\binom{N-m}{2},
    \\
    \gamma_3^{(3)}(m)
    &=32\sigma_4^{(3)}m\binom{m-1}{3}
      +24\sigma_4^{(3)}\binom m2(m-2)(N-m),
    \\
    \gamma_4^{(3)}(m)
    &=32\sigma_4^{(3)}\binom m4.
\end{align}

\subsubsection{Comparison with the large $N$ solution of Xu}

The Brownian clock of \cite{XuBrownianSpin} corresponds to
\begin{equation}
    \sigma_{q,{\rm Xu}}^{(3)}
    =\frac{N}{8q\,3^{q-1}\binom Nq}.
\end{equation}
With this normalization and fixed $m$,
\begin{equation}
    \gamma_1^{(3)}(m)\longrightarrow m,
    \qquad N\to\infty,
\end{equation}
while the remaining rates are suppressed.  The early-time master equation
becomes
\begin{equation}
    \partial_TP(m)
    =(m-q+1)P(m-q+1)-mP(m),
\end{equation}
so the elementary growth step and growth exponent are
\begin{equation}
    \Delta m=q-1,
    \qquad
    \lambda=q-1.
\end{equation}
In particular, $\lambda=1$ for $q=2$ and $\lambda=3$ for $q=4$.

The all-time large-$N$ solution of \cite{XuBrownianSpin} can be written
using
\begin{equation}
    x=\frac mN,
    \qquad
    x_{\rm eq}=\frac34,
    \qquad
    \lambda=q-1,
\end{equation}
\begin{equation}
    T_*
    =\frac1\lambda(\log N-\log\lambda),
\end{equation}
and
\begin{equation}
    p(x,T)
    =\frac{dz}{dx}
      \frac{e^{-z}z^{m_0/\lambda-1}}
           {\Gamma(m_0/\lambda)},
\end{equation}
where
\begin{equation}
    z(x,T)
    =\frac{x_{\rm eq}}{\lambda}
     e^{\lambda(T_*-T)}
     \left[
       -1+\left(1-\frac{x}{x_{\rm eq}}\right)^{-\lambda}
     \right].
\end{equation}
Changing from the Xu normalization to our Brownian convention only rescales
the clock,
\begin{equation}
    T_{\rm Xu}
    =\frac{\sigma_q^{(3)}}{\sigma_{q,{\rm Xu}}^{(3)}}T.
\end{equation}

\subsection{Brownian two-component Spin SYK}

Write $n=N_{\rm Spin}$.  The full Brownian Hamiltonian is
\begin{equation}
    H_{{\rm Spin},q}(t)
    =\sum_{1\le i_1<\cdots<i_q\le2n}
    J_{i_1\cdots i_q}(t)V_{i_1\cdots i_q},
\end{equation}
with
\begin{equation}
    \mathbb E[J_I(t)J_J(t')]
    =\sigma_q^{(2)}\delta_{IJ}\delta(t-t'),
    \qquad
    \sigma_q^{(2)}
    =\frac{(q-1)!{\cal J}_2}{(2n)^{q-1}}.
\end{equation}
The exact generator of the full model is
\begin{equation}
    G_{{\rm full},q}^{(k)}(Y)
    =\frac{\sigma_q^{(2)}}2
      \sum_{1\le i_1<\cdots<i_q\le2n}
      [W_I^{(k)},[W_I^{(k)},Y]],
\end{equation}
while for the genuine model
\begin{equation}
    G_{g,q}^{(k)}(Y)
    =\frac{\sigma_q^{(2)}}2
      \sum_{\substack{|S|=q\\
                      \boldsymbol\alpha\in\{x,y\}^q}}
      [W_{S,\boldsymbol\alpha}^{(k)},
       [W_{S,\boldsymbol\alpha}^{(k)},Y]].
\end{equation}

\subsubsection{Quadratic model}

For the genuine quadratic model,
\begin{equation}
    G_{g,2}^{(k)}(Y)
    =\frac{{\cal J}_2}{4n}
      \sum_{r<s}\sum_{\alpha,\beta=x,y}
      [W_{rs}^{\alpha\beta},[W_{rs}^{\alpha\beta},Y]].
\end{equation}
The full model also contains the $n$ same-site interactions proportional to
$\sigma_r^z$.  Since their signs drop out of the double commutator,
\begin{equation}
    G_{{\rm full},2}^{(k)}
    =G_{g,2}^{(k)}
     +\frac{{\cal J}_2}{4n}
      \sum_{r=1}^n[W_r^z,[W_r^z,\,\cdot\,]],
\end{equation}
where
\begin{equation}
    W_r^z=\sum_{a=1}^k(\sigma_r^z)^{(a)}.
\end{equation}

\subsubsection{Quartic model}

For the genuine quartic model,
\begin{equation}
    G_{g,4}^{(k)}(Y)
    =\frac{3{\cal J}_2}{8n^3}
      \sum_{r_1<r_2<r_3<r_4}
      \sum_{\alpha_1,\ldots,\alpha_4=x,y}
      [W_I^{(k)},[W_I^{(k)},Y]].
\end{equation}
The full quartic model has three classes of terms according to the number
$r$ of physical sites on which both constituent operators occur:
\begin{center}
\begin{tabular}{c|c|c}
$r$ & operator type & number of terms \\
\hline
$0$ & $\sigma^{\alpha_1}\sigma^{\alpha_2}
       \sigma^{\alpha_3}\sigma^{\alpha_4}$
    & $16\binom n4$ \\
$1$ & $\sigma^z\sigma^\alpha\sigma^\beta$
    & $4n\binom{n-1}{2}$ \\
$2$ & $\sigma^z\sigma^z$
    & $\binom n2$
\end{tabular}
\end{center}
Indeed,
\begin{equation}
    16\binom n4+4n\binom{n-1}{2}+\binom n2=\binom{2n}{4}.
\end{equation}
Thus
\begin{align}
    G_{{\rm full},4}^{(k)}
    =G_{g,4}^{(k)}
    +\frac{3{\cal J}_2}{8n^3}
    \bigg[&
      \sum_r\sum_{\substack{s<t\\s,t\neq r}}
      \sum_{\alpha,\beta=x,y}
      [W_{r;st}^{z;\alpha\beta},
       [W_{r;st}^{z;\alpha\beta},\,\cdot\,]]
      \nonumber\\
      &+\sum_{r<s}[W_{rs}^{zz},[W_{rs}^{zz},\,\cdot\,]]
    \bigg].
\end{align}
The two collision sectors contain $O(n^3)$ and $O(n^2)$ terms,
respectively, and are subleading to the $O(n^4)$ genuine sector at fixed
$q$ and large $n$.

\subsection{Operator-size dynamics in the two-component model}

The diagonal Pauli-string sector is invariant, but ordinary size alone is
not a closed variable.  For example, in the genuine $q=2$ model a single
$\sigma^z$ has $4(n-1)$ growth interactions, whereas a single $\sigma^x$
has only $2(n-1)$.

A minimal closed description keeps track of
\begin{equation}
    u=\#\{x\text{ or }y\text{ sites}\},
    \qquad
    v=\#\{z\text{ sites}\},
    \qquad
    m=u+v.
\end{equation}
For a genuine $q$-body interaction let
\begin{align}
    a&=\text{number of identity sites},\nonumber\\
    b&=\text{number of $z$ sites},\nonumber\\
    c&=\text{number of transverse sites with the same $x/y$ component},
    \nonumber\\
    d&=\text{number of transverse sites with the opposite $x/y$ component},
\end{align}
so that $a+b+c+d=q$.  The interaction anticommutes with the Pauli string
iff
\begin{equation}
    b+d\quad\text{is odd}.
\end{equation}
The number of interactions of type $(a,b,c,d)$ is
\begin{equation}
    {\cal N}_{abcd}(u,v)
    =2^{a+b}
      \binom{n-u-v}{a}\binom vb
      \frac{u!}{c!d!(u-c-d)!}.
\end{equation}
Multiplication by such an interaction sends
\begin{equation}
    (u,v)\longrightarrow
    {\cal T}_{abcd}(u,v)
    :=(u+a+b-c-d,\,v-b+d),
\end{equation}
and occurs at rate
\begin{equation}
    \Gamma_{abcd}(u,v)
    =4\sigma_q^{(2)}{\cal N}_{abcd}(u,v),
    \qquad b+d\ \text{odd}.
\end{equation}
Therefore
\begin{align}
    \partial_TP(u,v;T)
    =\sum_{\substack{a+b+c+d=q\\b+d\,\text{odd}}}
    \bigg[&
       \Gamma_{abcd}\!\left({\cal T}_{abcd}^{-1}(u,v)\right)
       P\!\left({\cal T}_{abcd}^{-1}(u,v);T\right)
       \nonumber\\
       &-\Gamma_{abcd}(u,v)P(u,v;T)
    \bigg],
\end{align}
with inadmissible predecessor states omitted.  The ordinary size changes by
\begin{equation}
    \Delta m=a-c,
\end{equation}
but the rate depends separately on $u$ and $v$.

\subsubsection{Exact quadratic master equation}

For $q=2$, writing $n_0=n-u-v$, the four transitions are
\begin{align}
    (u,v)&\longrightarrow(u,v+1),
    &\Gamma_+^{(\perp)}&=8\sigma_2^{(2)}un_0,
    \\
    (u,v)&\longrightarrow(u+2,v-1),
    &\Gamma_+^{(z)}&=16\sigma_2^{(2)}vn_0,
    \\
    (u,v)&\longrightarrow(u-2,v+1),
    &\Gamma_-^{(\perp)}&=4\sigma_2^{(2)}u(u-1),
    \\
    (u,v)&\longrightarrow(u,v-1),
    &\Gamma_-^{(z)}&=8\sigma_2^{(2)}uv.
\end{align}
The first two increase $m$ by one and the last two decrease it by one.  The
additional one-site $\sigma^z$ interactions in the full quadratic model
only rotate $\sigma^x\leftrightarrow\sigma^y$ and therefore leave $(u,v)$
unchanged.  The full and genuine models have the same coarse-grained
$(u,v)$ master equation although their full $k$-twirl generators differ.

\subsection{Large $n$ early-time growth in the two-component model}

At fixed $u,v$ and large $n$, a leading transition contains one occupied
site and $q-1$ identity sites.  If the occupied site is transverse,
\begin{equation}
    (u,v)\longrightarrow(u+q-2,v+1)
\end{equation}
with
\begin{equation}
    \Gamma_\perp(u,v)
    =4\sigma_q^{(2)}2^{q-1}u
      \binom{n-u-v}{q-1}
    \longrightarrow4{\cal J}_2u.
\end{equation}
If the occupied site is a $z$ site,
\begin{equation}
    (u,v)\longrightarrow(u+q,v-1)
\end{equation}
with
\begin{equation}
    \Gamma_z(u,v)
    =4\sigma_q^{(2)}2^qv
      \binom{n-u-v}{q-1}
    \longrightarrow8{\cal J}_2v.
\end{equation}
Both events increase $m$ by $q-1$, but with different rates.  The first
moments satisfy
\begin{equation}
    \frac{d}{dT}
    \begin{pmatrix}\langle u\rangle\\\langle v\rangle\end{pmatrix}
    =4{\cal J}_2
    \begin{pmatrix}
       q-2&2q\\
       1&-2
    \end{pmatrix}
    \begin{pmatrix}\langle u\rangle\\\langle v\rangle\end{pmatrix}.
\end{equation}
The dimensionless eigenvalues are
\begin{equation}
    \lambda_\pm^{(0)}
    =\frac{q-4\pm\sqrt{q(q+8)}}{2},
\end{equation}
so the dominant physical growth rate in this Brownian convention is
\begin{equation}
    \lambda_{\rm BDN}
    =2{\cal J}_2\left[q-4+\sqrt{q(q+8)}\right].
\end{equation}
In particular,
\begin{equation}
    \lambda_{\rm BDN}
    =\begin{cases}
       4{\cal J}_2(\sqrt5-1),&q=2,\\
       8\sqrt3\,{\cal J}_2,&q=4.
    \end{cases}
\end{equation}
An overall Brownian time rescaling changes these numerical values; the
structural statement is that the branching problem is two-component.

\subsection{The $\Gamma^3$ charge}

The two-component Brownian model retains
\begin{equation}
    \Gamma^3=\prod_{r=1}^n\sigma_r^z.
\end{equation}
For even $q$,
\begin{equation}
    [\Gamma^3,V_I]=0,
\end{equation}
so every replica charge commutes with the Brownian jump operators and the
semigroup cannot approach the unrestricted Haar twirl.  On a Pauli string
with $(u,v)$,
\begin{equation}
    \Gamma^3S(\Gamma^3)^{-1}=(-1)^uS.
\end{equation}
Moreover,
\begin{equation}
    \Delta u=a+b-c-d\equiv q\equiv0\pmod2,
\end{equation}
so the stochastic dynamics preserves $u\bmod2$.

\subsection{Comparison with the Brownian spin model of Xu}

For the Swingle--Winer model the identification with
\cite{XuBrownianSpin} is exact: the Brownian ensemble contains all
$3^q\binom Nq$ Pauli strings of size $q$.  Consequently the four-replica
Hamiltonian, the one-dimensional size master equation, and the large-$N$
solution above are precisely those of that work.

The Basu--Das--Nandy model differs because its local interaction alphabet is
restricted to $\{x,y\}$.  The diagonal Pauli-string sector remains closed,
but equal-size strings with different $z$ content have different rates.
Thus ordinary size $m$ must be replaced by $(u,v)$.  For the full quartic
model the collision terms modify the finite-$n$ master equation, but are
subleading at fixed $q$ and large $n$.


\section{Semigroup in Brownian SYK}

The Brownian SYK Hamiltonian is
\begin{equation}
    H(t)
    =\sum_{i_1<\cdots<i_q}
    J_{i_1\cdots i_q}(t)W_{i_1\cdots i_q},
    \qquad
    W_I=i^{q/2}\chi_{i_1}\cdots\chi_{i_q},
\end{equation}
with
\begin{equation}
    \mathbb E[J_I(t)]=0,
    \qquad
    \mathbb E[J_I(t)J_J(t')]
    =\sigma\delta_{IJ}\delta(t-t').
\end{equation}
For the normalization used above one may take
\begin{equation}
    \sigma=\frac{2^{q-1}q!}{q^2N^{q-1}}J,
\end{equation}
with $J$ now setting the Brownian clock.

The exact $k$-replica jump operators are
\begin{equation}
    X_I
    =\sqrt\sigma\sum_{a=1}^kW_I^{(a)}
    =\sqrt\sigma\,i^{q/2}
      \sum_{a=1}^k
      \chi_{i_1}^{(a)}\cdots\chi_{i_q}^{(a)},
\end{equation}
and
\begin{equation}
    G_{\rm eff}^{(k)}(Y)
    =\frac12\sum_I[X_I,[X_I,Y]].
\end{equation}
The doubled-space generator has exactly the algebraic form studied in
Sec.~4, but here the semigroup
\begin{equation}
    {\cal E}_{g;T}^{(k)}=e^{-TG_{\rm eff}^{(k)}}
\end{equation}
is exact at finite $N$.

\subsection{Time-dependent correlators}

Brownian independent increments also give a simple expression when the
replica evolution times are unequal.  Order
\begin{equation}
    0\le t_1\le t_2\le\cdots\le t_k.
\end{equation}
For $a=1,\ldots,k$, define the generator acting on the still-active
replicas $a,a+1,\ldots,k$,
\begin{equation}
    \widehat G_{[a:k]}
    :=\frac{\sigma}{2}\sum_I
      \left(
         \sum_{b=a}^k\widehat V_I^{(b)}
      \right)^2.
\end{equation}
Then
\begin{align}
    &\mathbb E_g\!\left[
       \bigotimes_{a=1}^k
       \left(U(t_a)\otimes U(t_a)^*\right)
    \right]
    \nonumber\\
    &\qquad=
    e^{-(t_k-t_{k-1})\widehat G_{[k:k]}}
    e^{-(t_{k-1}-t_{k-2})\widehat G_{[k-1:k]}}
    \cdots
    e^{-(t_2-t_1)\widehat G_{[2:k]}}
    e^{-t_1\widehat G_{[1:k]}}.
\end{align}
The order of the factors is important because the nested-subset generators
need not commute.  This formula is the natural starting point for
multi-time Brownian correlation functions.


\section{Approaching $k$-designs}

We now compare the Brownian moment channel
\begin{equation}
    {\cal E}_{g;T}^{(k)}=e^{-TG_{\rm eff}^{(k)}}
\end{equation}
with the $k$-fold Haar twirl
\begin{equation}
    {\cal T}_{\rm H}^{(k)}(Y)
    :=\int d\mu_{\rm Haar}(U)\,
      U^{\otimes k}YU^{\dagger\otimes k}.
\end{equation}

\subsection{Haar twirl}

By Schur--Weyl duality, the commutant of $U^{\otimes k}$ is spanned by
permutation operators $P_\sigma$, $\sigma\in S_k$.  We use
\begin{equation}
    P_\sigma|i_1\cdots i_k\rangle
    =|i_{\sigma(1)}\cdots i_{\sigma(k)}\rangle.
\end{equation}
For $d\ge k$ these operators are linearly independent.  The Gram matrix is
\begin{equation}
    G_{\tau,\sigma}
    :={\rm Tr}(P_\tau^\dagger P_\sigma)
    =d^{\#(\tau^{-1}\sigma)},
\end{equation}
where $\#(\pi)$ is the number of cycles of $\pi$.  Its inverse is expressed
through the unitary Weingarten function, giving
\begin{equation}
    {\cal T}_{\rm H}^{(k)}(Y)
    =\sum_{\sigma,\tau\in S_k}
      {\rm Wg}(\sigma\tau^{-1},d)
      {\rm Tr}(P_\tau^\dagger Y)P_\sigma.
\end{equation}
For $k=2$, with swap $S$,
\begin{equation}
    {\cal T}_{\rm H}^{(2)}(Y)
    =\frac{{\rm Tr}(Y)-d^{-1}{\rm Tr}(YS)}{d^2-1}\,I
     +\frac{{\rm Tr}(YS)-d^{-1}{\rm Tr}(Y)}{d^2-1}\,S.
\end{equation}

\subsection{Distance measures}

Several notions of distance are useful.

\paragraph{1. Hilbert--Schmidt distance.}
\begin{equation}
    \delta_k^{\rm HS}(T)
    :=\|{\cal E}_{g;T}^{(k)}-{\cal T}_{\rm H}^{(k)}\|_{2\to2},
\end{equation}
where
\begin{equation}
    \|\Phi\|_{2\to2}
    :=\sup_{Y\neq0}\frac{\|\Phi(Y)\|_2}{\|Y\|_2},
    \qquad
    \|Y\|_2=\sqrt{{\rm Tr}(Y^\dagger Y)}.
\end{equation}

\paragraph{2. Diamond distance.}
If $D=d^k$ is the input dimension,
\begin{equation}
    \delta_k^\diamond(T)
    :=\|{\cal E}_{g;T}^{(k)}-{\cal T}_{\rm H}^{(k)}\|_\diamond
    =\sup_{\rho\ge0,\ {\rm Tr}\rho=1}
      \left\|
       \bigl[({\cal E}_{g;T}^{(k)}-{\cal T}_{\rm H}^{(k)})
       \otimes{\rm id}_{D}\bigr](\rho)
      \right\|_1.
\end{equation}

\paragraph{3. Relative entropy distance.}
One may also consider
\begin{equation}
    \delta_k^{\rm RE}(T)
    :=\sup_\rho
      D\!\left(
        ({\cal E}_{g;T}^{(k)}\otimes{\rm id})(\rho)
        \,\middle\|\,
        ({\cal T}_{\rm H}^{(k)}\otimes{\rm id})(\rho)
      \right).
\end{equation}
This quantity requires the usual support conditions and is not controlled
by the spectral-gap argument alone.

\paragraph{4. Multiplicative or relative error.}
A standard CP-order notion is
\begin{equation}
    \epsilon_k(T)
    :=\inf\left\{
      \epsilon\ge0:
      (1-\epsilon){\cal T}_{\rm H}^{(k)}
      \preceq_{\rm CP}{\cal E}_{g;T}^{(k)}
      \preceq_{\rm CP}
      (1+\epsilon){\cal T}_{\rm H}^{(k)}
    \right\}.
\end{equation}
A priori none of these distances need tend to zero unless the steady-state
algebra of the Brownian dynamics coincides with the Haar invariant algebra.

\subsection{Steady states of Brownian dynamics}

The generator is
\begin{equation}
    G_{\rm eff}^{(k)}(Y)
    =\frac12\sum_I[X_I,[X_I,Y]].
\end{equation}
With the Hilbert--Schmidt inner product,
\begin{equation}
    \langle Y,G_{\rm eff}^{(k)}(Y)\rangle
    =\frac12\sum_I\|[X_I,Y]\|_2^2\ge0.
\end{equation}
Therefore
\begin{equation}
    \boxed{
    \ker G_{\rm eff}^{(k)}
    =\{Y:[X_I,Y]=0\ \text{for every }I\}.
    }
\end{equation}

Let $\mathfrak g$ be the one-copy Lie algebra generated by the
anti-Hermitian operators $iV_I$, and let ${\cal G}$ be the compact closure
of the corresponding unitary group.  Since
$X_I\propto\sum_aV_I^{(a)}$, the Lie algebra generated by the $X_I$ is the
diagonal $k$-copy representation of $\mathfrak g$.  Hence
\begin{equation}
    \ker G_{\rm eff}^{(k)}
    ={\rm Comm}\{U^{\otimes k}:U\in{\cal G}\}.
\end{equation}
The orthogonal projector onto this kernel is the subgroup Haar twirl
\begin{equation}
    {\cal T}_{\cal G}^{(k)}(Y)
    :=\int_{\cal G}d\mu_{\cal G}(U)\,
      U^{\otimes k}YU^{\dagger\otimes k}.
\end{equation}

Because $G_{\rm eff}^{(k)}$ is self-adjoint and positive semidefinite on the
finite-dimensional Hilbert--Schmidt space, if
\begin{equation}
    \Delta_k
    :=\min\{\lambda>0:\lambda\in{\rm spec}(G_{\rm eff}^{(k)})\}
\end{equation}
exists, then $\Delta_k>0$ and
\begin{equation}
    \left\|
       e^{-TG_{\rm eff}^{(k)}}-{\cal T}_{\cal G}^{(k)}
    \right\|_{2\to2}
    \le e^{-\Delta_kT}.
\end{equation}
Finite-dimensional norm equivalence then gives exponential convergence in
the diamond norm as well, with a dimension-dependent prefactor.

The Brownian dynamics converges to the full Haar twirl if and only if
\begin{equation}
    {\rm Comm}\{U^{\otimes k}:U\in{\cal G}\}
    ={\rm span}\{P_\sigma:\sigma\in S_k\}
\end{equation}
(for $d\ge k$).  A sufficient condition, valid for every $k$, is that the
one-copy dynamical Lie algebra contain $\mathfrak{su}(d)$.  For a fixed
finite $k$ this condition is sufficient but not necessary: a proper
subgroup can have the same $k$-fold commutant as the full unitary group.


\appendix
\section{Correlation functions and $k$-twirl maps}

Let ${\cal H}$ be the microscopic Hilbert space of dimension $d$.  On the
doubled Hilbert space define
\begin{equation}
    |1\rangle
    :=\frac1{\sqrt d}\sum_{i=1}^d|i\rangle_L|i\rangle_R,
    \qquad
    |A\rangle:=(A_L\otimes1_R)|1\rangle.
\end{equation}
Then
\begin{equation}
    \langle1|A_L\otimes B_R|1\rangle
    =\frac1d{\rm Tr}(AB^T).
\end{equation}

\paragraph{Permutation action and cycle states.}
On $k$ copies let $S_\pi$ permute the left factors.  For a $k$-cycle
$\pi$,
\begin{equation}
    \langle S_\pi1^{\otimes k}|
    \left(\prod_{a=1}^k(O_a)_L^{(a)}\right)
    |1^{\otimes k}\rangle
    =\frac1{d^k}
      {\rm Tr}\!\left(O_{\pi(k)}\cdots O_{\pi(1)}\right),
\end{equation}
with the ordering fixed by our convention for $S_\pi$.  In particular the
cyclic shift produces a cyclic trace.

\paragraph{Complexified evolution as a superoperator.}
Set the inverse temperature to one and define
\begin{equation}
    \rho=\frac{e^{-H}}{Z},
    \qquad Z={\rm Tr}(e^{-H}).
\end{equation}
For
\begin{equation}
    z=\theta+it,
    \qquad \theta\ge0,
\end{equation}
define
\begin{equation}
    U_z=e^{-zH},
    \qquad
    {\cal U}_z(O):=U_zOU_z^\dagger
    =e^{-zH}Oe^{-z^*H}.
\end{equation}
On the doubled space,
\begin{equation}
    \mathbf U_z
    =e^{-zH_L}\otimes e^{-z^*H_R^T}.
\end{equation}
For $\vec z=(z_1,\ldots,z_k)$ set
\begin{equation}
    \mathbf U_{\vec z}
    :=\bigotimes_{a=1}^k\mathbf U_{z_a}^{(a)}.
\end{equation}
The cycle identity gives
\begin{align}
    &{\rm Tr}\!\left[
       {\cal U}_{z_1}(O_1)\cdots{\cal U}_{z_k}(O_k)
    \right]
    \nonumber\\
    &\qquad=
    d^k
    \langle S_{\pi_0}1^{\otimes k}|
    \mathbf U_{\vec z}
    \left(\prod_{a=1}^k(O_a)_L^{(a)}\right)
    |1^{\otimes k}\rangle,
\end{align}
for the corresponding cyclic permutation $\pi_0$.  Thermal insertions are
obtained when
\begin{equation}
    \sum_{a=1}^k{\rm Re}\,z_a=\frac12,
\end{equation}
because the total Euclidean damping in the product is then $e^{-H}$.

\paragraph{Thermal one-point function.}
For any $t$,
\begin{equation}
    {\rm Tr}(\rho O)
    =\frac{
      \langle1|\mathbf U_{1/2+it}O_L|1\rangle
    }{
      \langle1|\mathbf U_{1/2+it}|1\rangle
    }
    =\frac{{\rm Tr}(e^{-H}O)}{Z}.
\end{equation}
The result is independent of $t$ by cyclicity.

\paragraph{Two-point function.}
For $0\le\tau\le1/2$ define
\begin{equation}
    G_{12}(t,\tau)
    :=\frac{
      {\rm Tr}\!\left[
        e^{-(1-\tau)H}O_1(t)e^{-\tau H}O_2
      \right]
    }{Z},
    \qquad
    O_1(t)=e^{iHt}O_1e^{-iHt}.
\end{equation}
It has the three-copy representation
\begin{equation}
    G_{12}(t,\tau)
    =\frac{
      \langle S_{\pi_0}1^{\otimes3}|
      \mathbf U_{(0,\tau,1/2-\tau)}
      \bigl((O_1(t))_L\otimes(O_2)_L\otimes1_L\bigr)
      |1^{\otimes3}\rangle
    }{
      \langle S_{\pi_0}1^{\otimes3}|
      \mathbf U_{(0,\tau,1/2-\tau)}
      |1^{\otimes3}\rangle
    }.
\end{equation}
More general thermal correlators can be represented similarly.  Identity
slots may be added when convenient so that the Euclidean intervals can be
encoded by nonnegative real parts of the $z_a$ satisfying the thermal-strip
constraint above.


\section{Conserved charges in Brownian SYK}

For the Brownian generator define
\begin{equation}
    {\rm Comm}\{X_I\}
    :=\{Y\in{\rm End}({\cal H}^{\otimes k}):[X_I,Y]=0\ \forall I\},
    \qquad
    X_I=\sqrt\sigma\sum_{a=1}^kV_I^{(a)}.
\end{equation}
As discussed in Sec.~8.3, this is the centralizer of the diagonal $k$-copy
representation of the one-copy dynamical group.

For even-$q$ SYK, each Hermitian interaction
\begin{equation}
    V_I=i^{q/2}\chi_{i_1}\cdots\chi_{i_q}
\end{equation}
is even under fermion parity.  On replica $a$ define
\begin{equation}
    Q^{(a)}
    :=(-i)^{N/2}\prod_{j=1}^N\chi_j^{(a)}.
\end{equation}
Then
\begin{equation}
    [Q^{(a)},X_I]=0
    \qquad\text{for all }a,I.
\end{equation}
Replica permutations also commute with every $X_I$, since they only relabel
the replica sum.  Hence
\begin{equation}
    {\rm Alg}\langle
       P_\sigma,Q^{(1)},\ldots,Q^{(k)}
    \rangle
    \subseteq {\rm Comm}\{X_I\}.
\end{equation}

If the one-copy algebra generated by the SYK interactions is the full even
matrix algebra
\begin{equation}
    {\rm End}({\cal H}_+)\oplus{\rm End}({\cal H}_-),
\end{equation}
then there are no additional tensor invariants.  In that case the commutant
is generated by replica permutations and replica parities and is spanned by
\begin{equation}
    P_\sigma\prod_{a=1}^k(Q^{(a)})^{\epsilon_a},
    \qquad
    \sigma\in S_k,
    \qquad
    \epsilon_a\in\{0,1\}.
\end{equation}
Equivalently, with parity projectors
\begin{equation}
    \Pi_{\vec s}
    :=\prod_{a=1}^k\frac{1+s_aQ^{(a)}}2,
    \qquad s_a=\pm1,
\end{equation}
the commutant is spanned by $\Pi_{\vec s}P_\sigma$.  If each replica is
restricted from the outset to a fixed parity sector, the $Q^{(a)}$ are
scalars and the commutant reduces to the permutation algebra, provided the
generated algebra is full on that sector.

    \bibliographystyle{unsrt}
    \bibliography{ref}
\end{document}